\begin{lemma}[Equipollence of cut and mix]
\label{lem:cut-mix-equipollence}
The calculus obtained from \(\LK\) by replacing \textup{\textsc{Cut}} by
\textup{\textsc{Mix}} derives exactly the \(\LK\)-derivable sequents, and likewise
for \(\LJ\).  Consequently, a derivation in the mix formulation with no
mix instance is a cut-free derivation, and eliminating mix eliminates
cut.
\textit{Status: complete proof.}
\end{lemma}
\begin{proof}
Cut is the mix instance \(m=n=1\), so reading each cut as a mix converts
a derivation with cuts into a derivation in the mix formulation with the
same endsequent.  Conversely, an instance of \textsc{Mix} with
multiplicities \(m\) and \(n\) is derived from the displayed primitive
rules: \(m-1\) applications of \textsc{CR} to the left premise contract
the succedent occurrences \(A^{m}\) to a single \(A\), \(n-1\)
applications of \textsc{CL} to the right premise contract \(A^{n}\) to a
single \(A\), and one \textsc{Cut} concludes
\(\Gamma,\Pi\Rightarrow\Delta,\Sigma\).  Replacing every mix instance,
uppermost first, by this derivation yields a derivation in the original
formulation with the same endsequent.  In \(\LJ\) the left premise
already has the single succedent \(A\), so only \textsc{CL} is used.
Hence the cut and mix formulations derive the same sequents, and a
mix-elimination procedure eliminates cut: first read every cut as the mix
instance \(m=n=1\), then eliminate the mixes.
\end{proof}

